\documentclass[../DoAn.tex]{subfiles}
\begin{document}
\section{Kết luận}
Đồ án đã xuất phát từ một vấn đề thực tế của học tập trực tuyến qua video: video bài giảng có khả năng tiếp cận rộng nhưng thường tạo ra trải nghiệm học thụ động, trong khi việc biến video thành bài học có tương tác lại đòi hỏi nhiều công sức thủ công từ giảng viên. Qua khảo sát các nhóm sản phẩm hiện có như nền tảng khóa học trực tuyến, công cụ tạo video tương tác và kho video công cộng, đồ án nhận thấy chưa nhóm sản phẩm nào kết hợp đồng thời ba khả năng: tự động sinh nội dung tương tác trực tiếp từ video, chấm điểm và phản hồi tự động, và hỗ trợ đa dạng loại hình tương tác trong cùng một luồng học. Trên cơ sở đó, đồ án đã xây dựng hệ thống Dyadia nhằm giảm công sức soạn bài cho giảng viên và tăng tính chủ động cho người học.

Kết quả chính của đồ án là một hệ thống web hoàn chỉnh gồm máy chủ ứng dụng richter, giao diện web heino và các thành phần hạ tầng đi kèm như PostgreSQL, FoundationDB, SeaweedFS, Speaches và Caddy. Hệ thống đã hiện thực được các nhóm chức năng quan trọng: quản lý tổ chức, khóa học và bài giảng; tải lên và phát video; phân tích video bằng trí tuệ nhân tạo; sinh các hoạt động tương tác theo phân đoạn nội dung; cho phép giảng viên xem trước, biên tập và xuất bản bài học; cho phép học viên học video kèm câu hỏi tương tác; chấm điểm, phản hồi và lưu kết quả; đồng thời cung cấp giao diện theo dõi kết quả học tập cho giảng viên. Các chức năng này tạo thành một quy trình khép kín từ video bài giảng ban đầu tới bài học tương tác đã được xuất bản và đánh giá.

So với các nền tảng học trực tuyến thông thường, Dyadia không chỉ quản lý khóa học và phát video mà còn gắn hoạt động kiểm tra trực tiếp vào dòng thời gian của video. So với các công cụ tạo video tương tác thủ công, hệ thống bổ sung khả năng tự động phiên âm, phân đoạn và sinh bài tập từ nội dung video. So với các kho video công cộng, hệ thống có cơ chế chấm điểm, phản hồi và phân tích kết quả học tập. Điểm khác biệt này không làm mất vai trò của giảng viên; ngược lại, hệ thống đặt giảng viên ở vai trò kiểm duyệt và tinh chỉnh kết quả sinh tự động trước khi đưa tới học viên.

Trong quá trình thực hiện, đồ án đã giải quyết được một số vấn đề kỹ thuật đáng chú ý. Thứ nhất là thiết kế quy trình nhiều giai đoạn để chuyển video thành bài học tương tác, gồm phiên âm lời thoại, phân đoạn nội dung và sinh hoạt động. Thứ hai là xây dựng mô hình hoạt động tương tác mở rộng được, hỗ trợ nhiều dạng bài như trắc nghiệm, điền khuyết, nghe hiểu, đọc hiểu và viết. Thứ ba là thiết kế hàng đợi tác vụ bền vững cho các xử lý dài, có checkpoint, cập nhật tiến độ và cơ chế bảo vệ quyền sở hữu worker. Thứ tư là kết hợp dữ liệu xem video với dữ liệu làm bài để tạo các thống kê học tập có ngữ cảnh, giúp giảng viên xác định đoạn nội dung khó và học viên cần hỗ trợ.

Về chất lượng triển khai, hệ thống đã được kiểm thử ở nhiều mức. Máy chủ ứng dụng có các ca kiểm thử đơn vị và tích hợp cho các luồng nghiệp vụ, dữ liệu và tác vụ nền. Giao diện web có các kịch bản kiểm thử đầu--cuối mô phỏng thao tác của người dùng trên trình duyệt. Kết quả kiểm thử trình bày ở Chương~\ref{chapter:Experiment} cho thấy các luồng chính của hệ thống vận hành đúng với thiết kế. Ngoài ra, hệ thống đã được đóng gói theo mô hình nhiều container để triển khai thực tế, trong đó Caddy đóng vai trò điểm vào, còn các dịch vụ nội bộ trao đổi qua mạng container.

Bên cạnh các kết quả đã đạt được, hệ thống vẫn còn một số hạn chế. Chất lượng nội dung sinh tự động phụ thuộc vào chất lượng âm thanh của video, độ chính xác của nhận dạng tiếng nói và khả năng của mô hình ngôn ngữ lớn. Với các bài giảng có tiếng ồn, nhiều người nói hoặc thuật ngữ chuyên ngành khó, kết quả phiên âm và câu hỏi sinh ra vẫn cần giảng viên rà soát kỹ. Hệ thống đã hỗ trợ nhiều loại tương tác, nhưng chưa bao phủ mọi hình thức đánh giá có thể có trong giáo dục, chẳng hạn bài tập lập trình, bài tập kéo thả phức tạp hoặc mô phỏng tương tác. Ngoài ra, phần phân tích học tập hiện tập trung vào các chỉ số trực tiếp như tiến độ xem, điểm số và tỷ lệ trả lời đúng; các mô hình dự báo sâu hơn về nguy cơ bỏ học hoặc cá nhân hóa lộ trình học vẫn chưa được phát triển.

Qua quá trình thực hiện, đồ án rút ra một số kinh nghiệm chính. Đối với các hệ thống có yếu tố trí tuệ nhân tạo, việc thiết kế dữ liệu trung gian và điểm kiểm soát của con người quan trọng không kém việc lựa chọn mô hình AI. Đối với các tác vụ dài, cần coi phục hồi lỗi, theo dõi tiến độ và tính nhất quán dữ liệu là yêu cầu thiết kế ngay từ đầu, thay vì xử lý bổ sung sau khi hệ thống đã hoàn thành. Đối với giao diện học tập, cần thiết kế đồng thời cho hai vai trò người dùng: học viên cần trải nghiệm đơn giản và ít gián đoạn, trong khi giảng viên cần đủ công cụ để kiểm soát chất lượng nội dung và theo dõi kết quả.

\section{Hướng phát triển}
Trong thời gian tới, hệ thống có thể tiếp tục được hoàn thiện theo một số hướng. Hướng đầu tiên là nâng cao chất lượng nội dung sinh tự động. Hệ thống có thể bổ sung cơ chế đánh giá chất lượng bản chép lời, phát hiện các đoạn có độ tin cậy thấp và gợi ý giảng viên kiểm tra lại trước khi sinh câu hỏi. Với bước sinh hoạt động tương tác, có thể bổ sung bộ tiêu chí kiểm tra tự động để phát hiện câu hỏi quá dễ, quá mơ hồ, trùng lặp hoặc không bám sát nội dung phân đoạn. Ngoài ra, hệ thống có thể cho phép giảng viên cung cấp thêm tài liệu tham chiếu như slide, giáo trình hoặc mục tiêu học tập để mô hình ngôn ngữ lớn sinh câu hỏi đúng trọng tâm hơn.

Hướng thứ hai là mở rộng hệ thống hoạt động tương tác. Các loại bài hiện có đã bao phủ những dạng phổ biến, nhưng hệ thống có thể tiếp tục bổ sung các dạng nâng cao như bài tập ghép cặp, sắp xếp thứ tự, kéo thả, câu hỏi theo tình huống, bài tập lập trình hoặc bài tập mô phỏng. Với mỗi loại mới, cần duy trì nguyên tắc đã áp dụng trong đồ án: có cấu hình dữ liệu rõ ràng, có cơ chế chấm điểm phù hợp, có giao diện biên tập cho giảng viên, giao diện làm bài cho học viên và giao diện xem lại kết quả.

Hướng thứ ba là phát triển sâu hơn phần phân tích học tập. Từ dữ liệu tiến độ xem, thời gian trả lời, số lần xem lại, điểm số và phản hồi theo từng phân đoạn, hệ thống có thể xây dựng các chỉ số nâng cao hơn như mức độ tập trung, xu hướng tiến bộ, nhóm câu hỏi gây nhầm lẫn hoặc nguy cơ học viên bỏ dở khóa học. Khi dữ liệu đủ lớn, có thể nghiên cứu mô hình gợi ý học lại phân đoạn, gợi ý bài tập bổ sung hoặc cảnh báo sớm cho giảng viên về những học viên cần hỗ trợ.

Hướng thứ tư là cải thiện khả năng vận hành ở quy mô lớn. Hệ thống hiện đã được tổ chức theo mô hình nhiều dịch vụ container, nhưng có thể tiếp tục bổ sung giám sát tập trung, thu thập log, cảnh báo lỗi, đo thời gian xử lý từng giai đoạn của pipeline và tự động điều chỉnh số lượng worker theo tải. Với các tác vụ xử lý video và nhận dạng tiếng nói, có thể nghiên cứu cơ chế phân bổ tài nguyên linh hoạt giữa CPU và GPU, cũng như hàng đợi ưu tiên cho các khóa học cần xử lý gấp.

Hướng thứ năm là hoàn thiện trải nghiệm cộng tác của giảng viên. Trong thực tế, một khóa học có thể do nhiều giảng viên hoặc trợ giảng cùng quản lý. Vì vậy, hệ thống có thể bổ sung cơ chế phân công người rà soát, bình luận trực tiếp trên từng phân đoạn hoặc từng câu hỏi, lưu lịch sử chỉnh sửa và so sánh các phiên bản nội dung. Những chức năng này sẽ giúp quá trình kiểm duyệt nội dung sinh tự động trở nên minh bạch và phù hợp hơn với môi trường đào tạo có nhiều người tham gia.

Cuối cùng, hệ thống có thể được đánh giá trên dữ liệu và người dùng thực tế với quy mô lớn hơn. Các thử nghiệm trong phạm vi đồ án chủ yếu nhằm chứng minh tính đúng đắn của luồng nghiệp vụ và khả năng vận hành kỹ thuật. Trong các bước tiếp theo, cần tổ chức thử nghiệm với nhiều giảng viên và học viên hơn để đo mức giảm thời gian soạn bài, mức cải thiện tỷ lệ hoàn thành bài học, mức độ hài lòng của người dùng và ảnh hưởng của tương tác trong video tới kết quả học tập. Đây sẽ là cơ sở quan trọng để tiếp tục hoàn thiện Dyadia thành một nền tảng học tập tương tác có thể ứng dụng trong môi trường đào tạo thực tế.
\end{document}
